\documentclass[sigconf]{acmart}
\AtBeginDocument{%
  \providecommand\BibTeX{{%
    Bib\TeX}}}

\setcopyright{acmlicensed}
\copyrightyear{2018}
\acmYear{2018}
\acmDOI{XXXXXXX.XXXXXXX}

\acmConference[SBQS25]{Brazilian Symposium on Software Quality}{Nov 04--07, 2025}{São José dos Campos, SP}

\settopmatter{printacmref=false}
\setcopyright{none}
\renewcommand\footnotetextcopyrightpermission[1]{}

\usepackage{multirow}
\usepackage{array}
\usepackage{graphicx,url}
\usepackage{float}
\usepackage{tabularx}
\usepackage[dvipsnames]{xcolor}
\definecolor{novocolor}{rgb}{0, 0.55, 0}
\newcommand{\novo}[1]{\textcolor{novocolor}{#1}}
\newcolumntype{C}{>{\centering\arraybackslash}X} % <-- modified
\usepackage[inline]{enumitem}
\usepackage{tcolorbox}
\usepackage[version=4]{mhchem} 

\usepackage{url}
\renewcommand{\url}[1]{#1}


\begin{document}

%%
%% The "title" command has an optional parameter,
%% allowing the author to define a "short title" to be used in page headers.
\title{Client-Side vs Server-Side Rendering: Impacts on Performance and User Experience in Web Applications}

\renewcommand{\shorttitle}{Client-Side vs Server-Side Rendering: Impacts on Performance and User Experience \newline in Web Applications }

%%
%% The "author" command and its associated commands are used to define
%% the authors and their affiliations.
%% Of note is the shared affiliation of the first two authors, and the
%% "authornote" and "authornotemark" commands
%% used to denote shared contribution to the research.

% \author{}
% \email{}
% \affiliation{%
%   \institution{\vspace{2cm}[{\em Omitted due to double-blind process}]\vspace{1.9cm}}
%   \city{}
%   \state{}
%   \country{}
% }

\author{Joey Clapton}
\affiliation{%
  \institution{Department of Software Engineering, Pontifical Catholic University of Minas Gerais (PUC Minas)}
  \city{Belo Horizonte}
  \state{Minas Gerais}
  \country{Brazil}
}
\email{joeyclapton42@gmail.com}

\author{Hugo de Paula}
\affiliation{%
  \institution{Department of Software Engineering, Pontifical Catholic University of Minas Gerais (PUC Minas)}
  \city{Belo Horizonte}
  \state{Minas Gerais}
  \country{Brazil}
}
\email{hugo@pucminas.br}

\author{João Pedro Oliveira Batisteli}
\affiliation{%
  \institution{Department of Software Engineering, Pontifical Catholic University of Minas Gerais (PUC Minas)}
  \city{Belo Horizonte}
  \state{Minas Gerais}
  \country{Brazil}
}
\email{joao.batisteli@sga.pucminas.br}


\author{Danilo Boechat Seufitelli}
\affiliation{%
  \institution{Federal Center for Technological Education of Minas Gerais (CEFET-MG)}
  \city{Belo Horizonte}
  \state{Minas Gerais}
  \country{Brazil}
}
\email{danilobs@cefetmg.br}


\author{Cleiton Tavares}
\orcid{1234-5678-9012}
\affiliation{%
  \institution{Department of Software Engineering, Pontifical Catholic University of Minas Gerais (PUC Minas)}
  \city{Belo Horizonte}
  \state{Minas Gerais}
  \country{Brazil}
}
\email{cleitontavares@pucminas.br}

%%
%% By default, the full list of authors will be used in the page
%% headers. Often, this list is too long, and will overlap
%% other information printed in the page headers. This command allows
%% the author to define a more concise list
%% of authors' names for this purpose.
\renewcommand{\shortauthors}{Clapton et al.}

%%
%% The abstract is a short summary of the work to be presented in the
%% article.

\renewcommand\abstractname{ABSTRACT}
\renewcommand\keywordsname{KEYWORDS}
\renewcommand\refname{REFERENCES}

\begin{abstract}
The increasing complexity of web applications necessitates careful selection of rendering strategies to ensure optimal performance and user experience. This study presents a comparative analysis of Client-Side Rendering (CSR) and Server-Side Rendering (SSR) techniques, evaluating their impact on web application quality based on Google's Core Web Vitals metrics (CWV).  In this work, photo gallery web applications were implemented using both rendering approaches to allow a homogeneous environment for comparable results.
Through controlled experiments across five data sizes (ranging from 498 kB to 26 MB), measurements were taken for key performance indicators, including First Contentful Paint (FCP), Speed Index (SI), Time to Interactive (TTI), Largest Contentful Paint (LCP), and JavaScript Heap Used. Statistical analysis using the Kruskal-Wallis test revealed significant performance differences ($p \leq$ 0.05) in all metrics. Results show that SSR achieves superior performance in TTI (mean reduction of 62.3\%) and LCP (mean reduction of 58.7\%), particularly under increased data loads, while CSR exhibits better FCP and SI performance for smaller datasets. The results suggest that the choice between CSR and SSR should consider the application's specific objectives and usage environment. 
These findings contribute to evidence-based decision-making frameworks for web application architecture, supporting software quality engineering practices in modern development environments.
\end{abstract}

%%
%% The code below is generated by the tool at http://dl.acm.org/ccs.cfm.
%% Please copy and paste the code instead of the example below.
%%
% \begin{CCSXML}
% <ccs2012>
%  <concept>
%   <concept_id>00000000.0000000.0000000</concept_id>
%   <concept_desc>Do Not Use This Code, Generate the Correct Terms for Your Paper</concept_desc>
%   <concept_significance>500</concept_significance>
%  </concept>
%  <concept>
%   <concept_id>00000000.00000000.00000000</concept_id>
%   <concept_desc>Do Not Use This Code, Generate the Correct Terms for Your Paper</concept_desc>
%   <concept_significance>300</concept_significance>
%  </concept>
%  <concept>
%   <concept_id>00000000.00000000.00000000</concept_id>
%   <concept_desc>Do Not Use This Code, Generate the Correct Terms for Your Paper</concept_desc>
%   <concept_significance>100</concept_significance>
%  </concept>
%  <concept>
%   <concept_id>00000000.00000000.00000000</concept_id>
%   <concept_desc>Do Not Use This Code, Generate the Correct Terms for Your Paper</concept_desc>
%   <concept_significance>100</concept_significance>
%  </concept>
% </ccs2012>
% \end{CCSXML}

% \ccsdesc[500]{Do Not Use This Code~Generate the Correct Terms for Your Paper}
% \ccsdesc[300]{Do Not Use This Code~Generate the Correct Terms for Your Paper}
% \ccsdesc{Do Not Use This Code~Generate the Correct Terms for Your Paper}
% \ccsdesc[100]{Do Not Use This Code~Generate the Correct Terms for Your Paper}

%%
%% Keywords. The author(s) should pick words that accurately describe
%% the work being presented. Separate the keywords with commas.
\keywords{Client-Side Rendering, Server-Side Rendering, Performance Analysis, User Experience, Web Applications}


% \received{20 February 2007}
% \received[revised]{12 March 2009}
% \received[accepted]{5 June 2009}

%%
%% This command processes the author and affiliation and title
%% information and builds the first part of the formatted document.
\maketitle

\section{Introduction}
\label{sec:introducao}

With the expansion of digital platforms, there is an increase in the volume of data generated and processed as a result of the digitization of daily activities and the growing need to collect and process large datasets~\cite{sqlBigData:2022}. In this context, the JavaScript Object Notation (JSON) format has emerged as a popular choice for representing and exchanging data in web applications due to its flexibility and ease of use. However, with its popularity, significant challenges have arisen, such as data parsing~\cite{misonJson:2017}, performance issues~\cite{dataInterchangeMobile:2011}, and concerns about the quality and efficiency of the application~\cite{pageExp:2022}. The increase in data volume implies higher computational costs. For this reason, it is crucial that developers understand the impact of their technological choices and coding practices on computational resource usage and the energy consumption of applications~\cite{pwaEnergy:2022}.


Previous studies compare data interaction formats in applications, addressing performance, computational costs, and environmental concerns. For instance, \citet{dataInterchangeMobile:2011} compare processing and energy consumption costs in mobile applications using XML, JSON, and Protocol Buffers.  \citet{dataConsumptionIot:2021} show that the choice of data encoding technology has a significant impact on performance and resource efficiency. Additionally, \citet{jsonXmlCaseStudy:2009} compare the performance of JSON and XML in a client-server environment, measuring transmission speed and computing resource usage. However, other factors such as the different processing approaches of data formats, both on the client and server sides, and their impact on energy consumption, performance, and user experience in data exchange for web applications must be considered.


Processing large volumes of JSON data has become a bottleneck in terms of performance and cost~\cite{misonJson:2017}, which highlights the importance of understanding and analyzing the cost of data interaction. Efficiently addressing this issue not only seeks to reduce the costs of handling data~\cite{misonJson:2017}, but also contributes to alleviating environmental concerns by reducing energy consumption and Carbon Dioxide (\ce{$CO_2$}) emissions. According to \citet{mvcEnergy:2023}, in the year 2020, the Information and Communication Technology (ICT) sector was responsible for 1.43 billion tons of \ce{$CO_2$}, representing 2.5\% of global greenhouse gas emissions.


In this scenario, the goal of this paper is to compare the efficiency and performance between Client-Side Rendering (CSR) and Server-Side Rendering (SSR) in web applications. To achieve this goal, the following specific objectives were proposed: (i) analyze performance aspects during data transactions; (ii) identify and compare the computational costs involved in data exchanges; and (iii) investigate and measure application performance and interactivity based on Google's Core Web Vitals (CWV) metrics.

%JP: adicionei uma breve descrição da metodologia adotada, sugiro uma revisão para verificar se está coerente. Também seria interessante adicionar mais detalhes da base de dados utilizada, como vi que faltam informações referentes a esse ponto no texto achei melhor não adicionar por agora.
To accomplish these objectives, a systematic methodology was adopted. A back-end application, implemented using Node.js, was made responsible for handling two parameters: data size and rendering type. A front-end application, developed in JavaScript, was designed to request the data, render the content, and display the results. Performance evaluation was automated through a custom script. This script leveraged the Lighthouse library for comprehensive auditing, \texttt{chrome-launcher} for initiating clean Chrome browser instances, and Puppeteer for collecting profiling and computational cost data.

The results demonstrate statistically significant differences between CSR and SSR across various metrics. It was observed that SSR rendering performed better in metrics such as Time to Interactive (TTI) and Largest Contentful Paint (LCP), while the CSR approach excelled in First Contentful Paint (FCP), and Speed Index (SI). Based on this data, it is expected that software engineers will be able to determine the most advantageous scenarios for each approach, taking into account response time, memory consumption, and CWV indicators. Thus, this study provides a guide to assist in choosing the most appropriate rendering method, contributing to the development of more efficient and user-friendly web applications that optimize performance and minimize computational costs.


This study is structured as follows: Section~\ref{sec:background} presents the theoretical background related to data interaction in web applications. Section~\ref{sec:related_work} discusses related work. Section~\ref{sec:materiais} details the materials and methods used in this research. Sections~\ref{sec:results} and~\ref{sec:discussion} present and discuss the results. Threats to validity are addressed in Section~\ref{sec:threats_to_validity}. Finally, Section~\ref{sec:conclusion} presents the conclusions.

\section{Background}
\label{sec:background}

This section presents the key concepts related to the study. Section \ref{sec:ssr} describes the SSR method. Section \ref{sec:csr} discusses CSR method. Section \ref{sec:performance_aspects} addresses software performance aspects, and Section \ref{sec:quality_indicators} covers quality indicators in web applications based on  Core Web Vitals (CWV) metrics.

%Fiz pequenas mudanças.
\subsection{Server-Side Rendering (SSR)}
\label{sec:ssr}

Server-Side Rendering (SSR) is a web application development methodology in which pages are rendered on the server and sent to the browser \cite{clientServerSide:2020}. When a request is made, the server processes it and returns the page already rendered in HTML to the user's browser.

% In applications that use SSR rendering is handled on the server side, allowing the page to be fully loaded when accessed by the user \citet{clientServerSide:2020}. This results in a faster user experience, especially on low-performance devices, since the server handles the more computationally expensive tasks. Additionally, SSR offers advantages in terms of search engine rankings, as it allows search engines to more easily index page content.

In applications that use SSR, the rendering process is handled on the server, enabling the page to be fully loaded when accessed by the user \citet{clientServerSide:2020}. This approach improves the user experience, particularly on low-performance devices, as the server performs the more computationally intensive tasks. Moreover, SSR provides advantages in search engine optimization (SEO), as it facilitates the indexing of page content by search engines.

\subsection{Client-Side Rendering (CSR)}
\label{sec:csr}

Client-Side Rendering is a web page rendering mechanism in which the rendering process takes place in the user's browser \cite{clientServerSide:2020}. The server sends a basic HTML structure with JavaScript responsible for dynamically rendering the content on the client side.

% According to \citet{clientServerSide:2020} CSR allows for more dynamic and interactive applications, as rendering logic and manipulation of the Document Object Model (DOM) occur directly in the user's browser. This can lead to a smoother user experience, especially in highly interactive applications such as Single Page Applications (SPAs), where navigation between pages does not require a full page reload.

CSR enables the development of more dynamic and interactive applications, as both the rendering logic and Document Object Model (DOM) manipulation are executed directly within the user's browser \cite{clientServerSide:2020}. This approach can deliver a smoother and more responsive user experience, particularly in highly interactive scenarios such as Single Page Applications (SPAs), where navigation between views occurs without the need for full page reloads.

\subsection{Software Performance Aspects}
\label{sec:performance_aspects}

A software performance aspect refers to a non-functional characteristic that describes the operational efficiency of a software system under various execution conditions~\cite{performanceAspects}. These aspects are essential to ensure that the software meets the expected quality and efficiency demanded by users. Key performance metrics include response time, which measures how long the system takes to respond to a request, and resource usage, which evaluates the memory and CPU consumption during application operations.

Performance aspects are critical to user experience, user retention, and the overall success of a web application. Pages that load slowly or respond poorly can frustrate users and cause them to abandon the site. Good performance is also important for Search Engine Optimization (SEO), as search engines tend to favor fast-loading websites~\cite{clientServerSide:2020}. As a result, optimizing these performance aspects has become a crucial part of developing and maintaining websites, requiring an integrated approach that considers both operational efficiency and end-user experience~\cite{pageExp:2022}.

\subsection{Quality Indicators in Web Applications Based on Core Web Vitals Metrics}
\label{sec:quality_indicators}

CWV is a set of performance metrics introduced by Google\footnote{\url{https://developers.google.com/web/tools/lighthouse}} to help optimize and evaluate the user experience of web pages, guiding improvements in the user experience. The CWV includes:

\textbf{Largest Contentful Paint (LCP)}: measures the loading time of the largest content element visible to the user.

\textbf{First Input Delay (FID)}: measures the time from the user's first interaction to the browser's response.

\textbf{Cumulative Layout Shift (CLS)}: evaluates layout shifts during page loading.

\textbf{First Contentful Paint (FCP)}: measures how quickly content is rendered to the screen, ensuring users see something useful sooner.

\textbf{Speed Index (SI)}: measures how quickly visible content is displayed during page loading, influencing user experience and search engine ranking.

\textbf{Time to Interactive (TTI)}: evaluates page usability. A faster TTI means users can interact with the page sooner.

These metrics are essential not only for ensuring a good user experience but also have significant implications for search engine rankings, as Google has started to consider user experience as a ranking factor for search results~\cite{pageExp:2022}.

\section{Related Work}
\label{sec:related_work}

The related works presented in this section discuss studies connected to the context of quality indicators in web applications based on Core Web Vitals (CWV) metrics, evaluation and comparison of performance aspects, and identification and comparison of computational costs in web applications.

\citet{parsingJson:2019} present  \textit{simdjson} tool, a high-performance solution for processing large volumes of data in JSON. The study emphasizes that, with the growing exchange of structured data in web applications, the performance of traditional parsers has become a bottleneck. \textit{Simdjson} outperforms reference analyzers such as RapidJSON by employing optimized techniques for detecting JSON structures. Thus, the study conceptually contributes to the present study by highlighting how the interpretation and manipulation of JSON data can directly impact overall application performance, particularly in contexts where different rendering strategies are compared.

\citet{webPerformanceOptimize:2023} analyzed an industrial case study in an agricultural technology company aiming to improve the performance of a large-scale web dashboard. The authors developed a performance engineering plan using a benchmarking tool based on Google Lighthouse to assess the impact on web performance metrics. The results demonstrated a significantly faster browsing experience, with pages loading 98.37\% faster on desktop devices and 97.56\% faster on mobile. Content rendering also became more efficient, with improvements of 48.25\% and 19.85\% in rendering speed on desktop and mobile devices, respectively. This study is relevant to the current research as it provides data on the impact of different interventions affecting performance metrics using Google Lighthouse. However, there remains a gap for analyzing performance during data transactions between JSON and HTML formats using the same tool.

\citet{webAutoscaling:2022} explored mechanisms to ensure the resilience and performance of web applications under high-demand scenarios and potential failures. The study investigated different scaling approaches for applications and databases, such as horizontal and vertical scaling, database replication, and sharding, focusing on achieving fault tolerance and autoscaling capabilities. The authors emphasize the importance of making applications stateless by using external solutions for session storage and user files to facilitate horizontal scaling. The study concludes that to build a fault-tolerant web application, scalability is essential, and autoscaling can significantly reduce manual configuration time and improve responsiveness to traffic spikes. The research offers insights into how different scaling strategies can influence the performance and resilience of web applications, critical aspects also considered in the comparison between CSR and SSR.

\citet{clientServerSide:2020} investigated the performance of web applications built using CSR and SSR rendering techniques. The study developed a web application named Show Up using both approaches and compared their performance based on FCP, Speed Index, TTI, and FMP metrics. The results showed that SSR rendering outperformed CSR across all evaluated metrics, delivering a faster and smoother user experience. In addition, SSR scored higher on Google audit metrics for performance, accessibility, and best practices. The authors provide a discussion on different rendering strategies and their implications for application performance. This enables further investigation into other contributing factors, such as the impact of JSON and HTML data formats on user experience and web application efficiency, analyzed in the present study.

\section{Design Study}
\label{sec:materiais}

This study conducts a quantitative experimental research by investigating and comparing the performance, computational cost, and CWV metrics analysis between SSR and CSR for data presentation in web applications. The following subsections describe the development environment, which consists of a web application. Subsequently, the materials and methods used in the experiment are detailed, including performance measurement tools, computational cost metrics, and the different data sizes analyzed.

\subsection{Study Steps}

The study was divided into three steps, which are described as follows.

% This section presents the three main stages of development for this study.

In the \textit{first step}, the data were defined and created in JSON format. The datasets were generated in five different sizes: 498 kB, 2 MB, 7.9 MB, 16.8 MB, and 26 MB. Each dataset was composed by combining 13 different images, resulting in: 370 images (498 kB), 1,500 images (2 MB), 6,000 images (7.9 MB), 12,000 images (16.8 MB), and 19,000 images (26 MB).
These sizes were selected based on the study by~\citet{jsonXmlCaseStudy:2009}, which analyzed the performance of data interchange formats in various scenarios. For the experiment, a data structure was designed to simulate a photo gallery system using JSON (Figure~\ref{fig:jsonStructure}). The choice of this data structure aims to represent a real-world use case where user information is frequently exchanged between systems.

\begin{figure}[h]
\centering
\includegraphics[width=.45\textwidth]{img/03-json-structure.png}
\caption{JSON data structure used in the application}
\label{fig:jsonStructure}
\end{figure}

Furthermore, a back-end application was developed using the NodeJS platform. An Application Programming Interface (API) was created to handle two parameters: data size and rendering type (SSR or CSR). Based on these parameters, the API returns either the HTML rendered from the JSON or the raw JSON itself to the client. The data is stored locally.

The \textit{second step} %involved developing a front-end application in JavaScript. This application (Figure~\ref{fig:exampleFig2}) allows users to select both the size and format of the data to be requested. It then sends the request to the back-end and displays the response, which may either be the fully rendered HTML from the server or the raw JSON used to dynamically build the web page.
consisted of developing a front-end application in JavaScript. As illustrated in Figure~\ref{fig:exampleFig2}, this application enables users to select both the size and format of the data to be requested. Once the request is sent to the back-end, the application displays the response, which can either be the fully rendered HTML returned by the server or the raw JSON used to dynamically build the web page.

\begin{figure}[h]
\centering
\includegraphics[width=.5\textwidth]{img/04-csr-ssr.png}
\caption{Application operation flow}
\label{fig:exampleFig2}
\end{figure}

Subsequently, a JavaScript script was developed to evaluate application performance. %The script uses the Lighthouse library to automatically calculate key performance metrics. Additionally, the \texttt{chrome-launcher} library was employed to launch a clean instance of the Chrome browser for performance testing. The Puppeteer tool\footnote{\url{https://pptr.dev/}} was also used to collect profiling metrics and computational cost data. Each configuration (JSON and HTML) was tested 10 times, and the average execution time was computed to yield more accurate and representative results. The following metrics, were analyzed in this study:
Leveraging the Lighthouse library, the script automatically computes key performance metrics. To ensure consistent testing conditions, the \texttt{chrome-launcher} library was used to start a clean instance of the Chrome browser. In addition, the Puppeteer tool\footnote{\url{https://pptr.dev/}} was used to collect profiling data and computational cost metrics. Each configuration (JSON and HTML) was tested 10 times, with the average execution time calculated to provide more accurate and representative results. The following metrics were analyzed in this study:

%%=============================
% Pra deixar o texto auto-contido, eu add uma breve explicação. Caso tenha algo errado, favor verificar Joey
% Na real percebi que ja tem explicação delas, exceto a última.
\begin{enumerate}
    \item \textbf{First Contentful Paint}%: Measures the time until the first piece of content (e.g. text or image) renders on the screen.
    \item \textbf{Speed Index}%: Measures how quickly the content within the visible area (viewport) is visually populated.
    \item \textbf{Time to Interactive}%: Measures the time for a page to become fully responsive to user interactions.
    \item \textbf{Largest Contentful Paint}%: Measures the render time of the largest content element (image or text block) within the visible area.
    \item \textbf{\textit{JS Heap} Size}%: Measures memory usage by JavaScript, indicating the application's computational cost.
\end{enumerate}

In the \textit{third step}, the Kruskal-Wallis statistical test\footnote{Kruskal-Wallis is a non-parametric test used for comparing three or more independent samples.} was applied to validate the significance of the results across the different data formats tested. Core Web Vitals metrics were evaluated to understand the impact of data formats on user experience.



\subsection{Development Environment}

All experiments were conducted in a controlled environment to ensure consistency and reproducibility of results. The system's development and execution were performed on a MacBook Air M2, featuring an 8-core \textit{3.5 GHz} CPU, running macOS Sequoia 15.0, equipped with 16GB of RAM and a 256GB SSD.

% \section{Development}
% \label{sec:development}

% This section presents the operation of the application and the measurement tools used in the tests.

\subsection{Photo Application}
Two photo gallery applications were developed: one focused on server-side rendering, and the other on client-side rendering, along with a selection page where the user can choose which rendering method will be displayed. Both applications were built using the same HTML structure, the same CSS styles, and without inserting any scripts, in order to prevent side effects on test results. Figure~\ref{fig:jsonStructure1} shows the initial selection page, while Figure~\ref{fig:jsonStructure2} shows the application running with client-side rendering using the smallest data size.

% \begin{figure}[h]
% \centering
% \begin{minipage}{0.47\textwidth}
%     \centering
%     \includegraphics[width=\textwidth]{img/app02.png}
%     \caption{Data size and rendering type selection}
%     \label{fig:jsonStructure1}
% \end{minipage}\hfill
% \begin{minipage}{0.47\textwidth}
%     \centering
%     \includegraphics[width=\textwidth]{img/app-client.png}
%     \caption{Application using extra small data and CSR rendering}
%     \label{fig:jsonStructure2}
% \end{minipage}
% \end{figure}




\begin{figure}[h]
\centering
    \includegraphics[width=.45\textwidth]{img/app02.png}
    \caption{Data size and rendering type selection}
    \label{fig:jsonStructure1}
\end{figure}

\begin{figure}[h]
\centering
    \includegraphics[width=.45\textwidth]{img/app-client.png}
    \caption{Application using extra small data and CSR}
    \label{fig:jsonStructure2}
\end{figure}

\begin{table*}[htbp]
\centering
\scriptsize
\caption{Analyzed Metrics}
\begin{tabular}{|l|l|c|c|c|c|c|c|c|c|c|c|}
\hline
\textbf{Metric} & \textbf{Measures} & \textbf{Client Small} & \textbf{Server Small} & \textbf{Client XS} & \textbf{Server XS} & \textbf{Client M} & \textbf{Server M} & \textbf{Client L} & \textbf{Server L} & \textbf{Client XL} & \textbf{Server XL} \\
\hline
\multirow{3}{*}{\textbf{LCP (ms)}} &
Average & 14,716 & 4,520 & 17,382 & 5,270 & 28,841 & 10,456 & 49,067 & 20,245 & 68,644 & 31,369 \\
& Median & 14,716 & 3,837 & 17,779 & 2,180 & 28,839 & 10,456 & 49,051 & 20,240 & 69,325 & 30,342 \\
& Std. Dev. & 1 & 2,048 & 920 & 6,180 & 5 & 1 & 31 & 8 & 2,185 & 3,078 \\
\hline
\multirow{3}{*}{\textbf{FCP (ms)}} &
Average & 608 & 2,980 & 614 & 1,241 & 608 & 9,668 & 608 & 19,424 & 608 & 29,548 \\
& Median & 608 & 2,979 & 609 & 1,240 & 608 & 9,668 & 608 & 19,433 & 608 & 29,542 \\
& Std. Dev. & 0.6 & 3.6 & 8.8 & 3.4 & 1.3 & 3.2 & 0.9 & 31 & 0.6 & 19 \\
\hline
\multirow{3}{*}{\textbf{SI (ms)}} &
Average & 608 & 2,980 & 614 & 1,241 & 792 & 9,668 & 1,250 & 19,424 & 1,796 & 29,548 \\
& Median & 608 & 2,979 & 609 & 1,240 & 785 & 9,668 & 1,249 & 19,433 & 1,793 & 29,542 \\
& Std. Dev. & 0.6 & 3.6 & 8.8 & 3.4 & 23 & 3.2 & 27 & 31 & 27 & 19 \\
\hline
\multirow{3}{*}{\textbf{TTI (ms)}} &
Average & 15,646 & 4,593 & 17,794 & 5,370 & 29,816 & 10,509 & 50,035 & 20,493 & 69,452 & 32,220 \\
& Median & 15,646 & 3,890 & 18,154 & 2,232 & 29,814 & 10,509 & 50,019 & 20,481 & 70,199 & 31,120 \\
& Std. Dev. & 1 & 2,108 & 832 & 6,276 & 5 & 1 & 32 & 26 & 2,404 & 2,867 \\
\hline
\multirow{3}{*}{\textbf{\textit{JS Heap} Used (KB)}} &
Average & 202,997 & 205,631 & 120,867 & 108,542 & 440,358 & 436,163 & 525,017 & 491,820 & 613,318 & 560,978 \\
& Median & 205,728 & 207,186 & 121,014 & 104,708 & 446,453 & 443,045 & 528,830 & 485,718 & 612,958 & 552,913 \\
& Std. Dev. & 12,000 & 10,105 & 3,696 & 10,222 & 26,620 & 15,233 & 31,953 & 27,080 & 32,861 & 33,344 \\
\hline
\end{tabular}
\vspace{2pt}
\begin{flushleft}
\footnotesize \textbf{*} Client = Client Side, Server = Server Side, XS = Extra Small, M = Medium, L = Large, XL = Extra Large, Std. Dev. = Standard Deviation.
\end{flushleft}
\label{tab:analyzed_metrics}
\end{table*}


\begin{figure*}[h]
    \centering
    \includegraphics[width=\textwidth, keepaspectratio]{img/lcp.png}
    \caption{LCP values}
    \label{fig:lcp}
\end{figure*}

\begin{figure*}[h]
    \centering
    \includegraphics[width=\textwidth, keepaspectratio]{img/fcp.png}
    \caption{FCP values}
    \label{fig:fcp}
\end{figure*}

\section{Results}
\label{sec:results}

After conducting the experiment, the average, median, and standard deviation values were calculated for each metric and rendering method. It was observed that, for smaller data sizes, the metrics showed lower and more consistent values, especially on the server side. On the other hand, as the data size increased, the standard deviation also increased in some measurements, indicating greater variability in the results. Table~\ref{tab:analyzed_metrics} presents the computed results.

\novo{Nesta versão estendida, o design experimental foi ampliado de 10 condições originais (5 tamanhos $\times$ 2 abordagens de renderização) para 30 condições, mediante a introdução de dois fatores adicionais: minificação de JSON e compressão HTTP via gzip. A minificação foi aplicada exclusivamente às respostas CSR (formato JSON), uma vez que o SSR retorna HTML pré-renderizado, que não foi submetido a minificação neste estudo. O design fatorial resultante compreende seis condições de tratamento: CSR (baseline), CSR + Minificação, CSR + Gzip, CSR + Minificação + Gzip, SSR (baseline) e SSR + Gzip. Cada condição foi executada 10 vezes por tamanho de conjunto de dados, totalizando 300 observações. As Tabelas~\ref{tab:lcp-medians} e~\ref{tab:tti-medians} apresentam os valores medianos de LCP e TTI para todas as 30 condições.}

% Tabela: Medianas de LCP (segundos) por condição e tamanho de conjunto de dados
\begin{table*}[htbp]
\centering
\footnotesize
\caption{\novo{Medianas de LCP (segundos) por condição e tamanho de conjunto de dados}}
\begin{tabular}{|l|c|c|c|c|c|}
\hline
\textbf{Condição} & \textbf{Extra-Small} & \textbf{Small} & \textbf{Medium} & \textbf{Large} & \textbf{Extra-Large} \\
\hline
CSR (baseline)               & 4,21  & 11,64 & 40,62 & 83,85  & 127,39 \\
CSR + Minificação            & 3,08  &  6,46 & 20,05 & 40,32  &  60,59 \\
CSR + Gzip                   & 2,26  &  2,48 &  3,96 &  6,35  &   8,58 \\
CSR + Minificação + Gzip     & 2,11  &  2,41 &  3,42 &  5,25  &   6,95 \\
SSR (baseline)               & 1,53  &  3,19 &  9,65 & 19,57  &  29,35 \\
\textbf{SSR + Gzip}          & \textbf{0,93} & \textbf{0,94} & \textbf{1,10} & \textbf{1,28} & \textbf{1,46} \\
\hline
\end{tabular}
\vspace{2pt}
\begin{flushleft}
\footnotesize
\textbf{*} Em negrito: melhor resultado por tamanho.
Tamanhos de conjunto de dados: Extra-Small $\approx$ 498~kB; Small $\approx$ 2~MB; Medium $\approx$ 7,9~MB; Large $\approx$ 16,8~MB; Extra-Large $\approx$ 26~MB.
\end{flushleft}
\label{tab:lcp-medians}
\end{table*}


% Tabela: Medianas de TTI (segundos) por condição e tamanho de conjunto de dados
\begin{table*}[htbp]
\centering
\footnotesize
\caption{\novo{Medianas de TTI (segundos) por condição e tamanho de conjunto de dados}}
\begin{tabular}{|l|c|c|c|c|c|}
\hline
\textbf{Condição} & \textbf{Extra-Small} & \textbf{Small} & \textbf{Medium} & \textbf{Large} & \textbf{Extra-Large} \\
\hline
CSR (baseline)               & 4,24  & 11,66 & 40,64 & 83,87  & 127,41 \\
CSR + Minificação            & 3,10  &  6,49 & 20,08 & 40,35  &  60,62 \\
CSR + Gzip                   & 2,27  &  2,50 &  3,99 &  6,37  &   8,60 \\
CSR + Minificação + Gzip     & 2,14  &  2,44 &  3,44 &  5,27  &   6,97 \\
SSR (baseline)               & 1,53  &  3,19 &  9,65 & 19,57  &  29,35 \\
\textbf{SSR + Gzip}          & \textbf{0,93} & \textbf{0,94} & \textbf{1,10} & \textbf{1,28} & \textbf{1,46} \\
\hline
\end{tabular}
\vspace{2pt}
\begin{flushleft}
\footnotesize
\textbf{*} Em negrito: melhor resultado por tamanho. Valores em segundos.
\end{flushleft}
\label{tab:tti-medians}
\end{table*}



\paragraph{\textbf{Largest Contentful Paint}}
The LCP metric measures the loading time of the largest element on the page, which may be an image, text section, or video. A value up to 2.5s indicates a good user experience; between 2.5s and 4.0s indicates room for improvement; and above 4.0s indicates a poor experience. Figure~\ref{fig:lcp} shows the boxplot with the results. SSR rendering achieved better performance across all data sizes compared to CSR, which performed worse, especially with larger datasets.

\novo{A habilitação da compressão gzip produziu a maior melhoria em LCP em ambas as abordagens de renderização. Para o CSR, o gzip reduziu o LCP mediano em 46,3\% para payloads Extra-Small e em até 93,3\% para Extra-Large (de 127,39~s para 8,58~s). Para o SSR, o gzip reduziu o LCP mediano de 29,35~s para 1,46~s no tamanho Extra-Large (redução de 95,0\%), mantendo valores dentro do limiar ``bom'' ($\leq 2{,}5$~s) em todos os tamanhos de dados.

A minificação de JSON, aplicada às respostas CSR sem gzip, também reduziu o LCP de forma significativa: de 127,39~s para 60,59~s no conjunto Extra-Large (redução de 52,4\%) e de 4,21~s para 3,08~s no Extra-Small (redução de 26,8\%). Contudo, a melhoria obtida apenas com minificação foi consideravelmente menor do que a proporcionada pelo gzip. Quando ambas as técnicas foram combinadas no CSR, o LCP foi reduzido para 6,95~s no Extra-Large, representando apenas uma melhoria adicional de 19,0\% em relação ao CSR + Gzip isolado (8,58~s), indicando retorno marginal decrescente da minificação quando o gzip já está ativo.

Em todas as condições e tamanhos, o SSR + Gzip alcançou os menores valores de LCP. Mesmo comparado com a melhor configuração CSR (CSR + Minificação + Gzip), o SSR + Gzip permaneceu mais rápido por um fator crescente com o tamanho do payload: 2,3$\times$ para Extra-Small, 2,6$\times$ para Small, 3,1$\times$ para Medium, 4,1$\times$ para Large e 4,8$\times$ para Extra-Large. Esse padrão é consistente com a análise de~\citet{ollila:2021}, que demonstrou que as diferenças de desempenho entre estratégias de renderização são amplificadas conforme o tamanho dos dados de entrada aumenta.}


\paragraph{\textbf{First Contentful Paint}}
The FCP metric refers to the time between when the user navigates to the page and when the first piece of content is rendered on the screen. A good experience occurs when this value is up to 1.8s, while values above 3s are considered poor. Figure~\ref{fig:fcp} presents the boxplot with the results. For smaller datasets, both server-side and client-side rendering provide fast user experiences. As data volume increases, SSR shows a significant increase in loading time, while CSR maintains consistently good performance compared to SSR.

\novo{Um resultado arquiteturalmente relevante surgiu na análise do FCP para o CSR: a métrica permaneceu efetivamente constante em aproximadamente 1,58~s, independentemente da aplicação de minificação ou gzip (por exemplo, no Extra-Small, todas as quatro condições CSR produziram FCP médio de $1{,}580 \pm 0{,}001$~s). Essa invariância ocorre porque, no CSR, o FCP é disparado quando o navegador renderiza o HTML inicial contendo o bootstrap JavaScript — um artefato de tamanho fixo carregado antes da busca do payload de dados. Portanto, nem a compressão gzip nem a minificação do JSON afetam o instante em que o navegador pinta o primeiro pixel. Esse achado indica que o FCP \emph{não} é uma métrica adequada para avaliar o impacto de técnicas de otimização de transferência de dados em aplicações CSR; LCP e TTI são métricas mais discriminantes para esse fim.

Para o SSR, o comportamento do FCP foi substancialmente diferente e fortemente dependente do gzip. Sem compressão, o FCP escalou proporcionalmente com o tamanho do payload, variando de 1,23~s (Extra-Small) a 29,28~s (Extra-Large), pois o navegador não consegue renderizar nenhum conteúdo até que uma parcela suficiente do documento HTML pré-renderizado seja recebida. Com o gzip habilitado, o FCP do SSR caiu para valores entre 0,63~s e 1,40~s em todos os tamanhos — todos dentro do limiar ``bom'' do Lighthouse —, representando reduções de 38,9\% a 95,2\% em relação ao baseline sem compressão. Conforme observado por~\citet{verma:2024}, o SSR oferece tempos de carregamento inicial mais rápidos quando a resposta do servidor é entregue de forma eficiente.}


\paragraph{\textbf{Speed Index}}
The SI metric measures the speed at which content is displayed during page loading. Lighthouse classifies times up to 3.4s as fast, between 3.4s and 5.8s as moderate, and above 5.8s as slow. Figure~\ref{fig:si} shows the boxplot with the results. For smaller data sizes, SI followed the same trend as the LCP and FCP metrics, with both rendering methods showing comparable loading times. However, with larger data volumes, there was a significant performance degradation in server-side rendering, resulting in worse outcomes.

\begin{figure*}[h]
    \centering
    \includegraphics[width=\textwidth, keepaspectratio]{img/si.png}
    \caption{SI values}
    \label{fig:si}
\end{figure*}

\novo{Para o CSR, o SI permaneceu constante em aproximadamente 1,58~s independentemente do nível de otimização, uma vez que o conteúdo visível começa a ser renderizado imediatamente quando o shell JavaScript é pintado. Para o SSR, o SI exibiu a mesma tendência do FCP: sem gzip, escalou com o tamanho do payload até 29,28~s no Extra-Large; com gzip, foi reduzido para valores abaixo de 1,5~s em todos os tamanhos. Esses resultados indicam que o gzip é condição necessária para que o SSR alcance desempenho competitivo de SI em cenários de payload elevado.}


\paragraph{\textbf{Time to Interactive}}
The TTI metric measures how long it takes for the page to become fully interactive. Lighthouse considers times up to 3.4s as fast, between 3.4s and 5.8s as moderate, and above 5.8s as slow. Figure~\ref{fig:tti} shows the boxplot with the results. SSR showed shorter loading times compared to CSR across all data sizes.

\begin{figure*}[h]
    \centering
    \includegraphics[width=\textwidth, keepaspectratio]{img/tti.png}
    \caption{TTI values}
    \label{fig:tti}
\end{figure*}

\novo{Diferentemente do FCP e do SI, o TTI para CSR foi substancialmente afetado tanto pelo gzip quanto pela minificação. Como o CSR exige o recebimento e a análise completa do payload de dados antes de a página tornar-se interativa, o TTI rastreia o LCP de perto (diferindo em aproximadamente 20--30~ms), conforme evidenciado na Tabela~\ref{tab:tti-medians}. O SSR + Gzip produziu consistentemente valores de TTI iguais ou inferiores a 1,5~s em todos os tamanhos de conjunto de dados — classificados como ``rápido'' pelo Lighthouse. O CSR sem otimização atingiu TTI de até 127,41~s no Extra-Large, caracterizando uma experiência de usuário extremamente ruim por qualquer critério~\cite{sikder:2016}.}


\paragraph{\textbf{JS Heap Size}}
The \textit{JS Heap Size} metric refers to the memory space allocated for storing JavaScript objects. Although there is no standardized quality threshold, it was observed that CSR and SSR had similar memory usage for smaller data sizes. However, memory consumption increased significantly for CSR with larger datasets. Figure~\ref{fig:jsheap} shows the boxplot with the results.

\begin{figure*}[h]
    \centering
    \includegraphics[width=\textwidth, keepaspectratio]{img/jsheap.png}
    \caption{\textit{JS Heap} memory usage}
    \label{fig:jsheap}
\end{figure*}


\novo{\paragraph{\textbf{Lighthouse Performance Score}}
O Lighthouse Performance Score é uma métrica composta que pondera FCP, SI, LCP, TTI e Total Blocking Time em uma escala de 0 a 100. A Tabela~\ref{tab:lighthouse-score} apresenta as medianas de Score por condição e tamanho.

% Tabela: Medianas do Lighthouse Performance Score por condição e tamanho
\begin{table*}[htbp]
\centering
\footnotesize
\caption{\novo{Medianas do Lighthouse Performance Score (0--100) por condição e tamanho de conjunto de dados}}
\begin{tabular}{|l|c|c|c|c|c|}
\hline
\textbf{Condição} & \textbf{Extra-Small} & \textbf{Small} & \textbf{Medium} & \textbf{Large} & \textbf{Extra-Large} \\
\hline
CSR (baseline)              & 85  & 74 & 74 & 74 & 74 \\
CSR + Minificação           & 88  & 77 & 77 & 77 & 77 \\
CSR + Gzip                  & 87  & 97 & 87 & 87 & 87 \\
\textbf{CSR + Min. + Gzip}  & \textbf{98} & \textbf{97} & \textbf{98} & \textbf{98} & \textbf{98} \\
SSR (baseline)              & 77  & 65 & $<30$ & $<15$ & $<10$ \\
SSR + Gzip                  & 77  & 77 & 75 & 74 & 74 \\
\hline
\end{tabular}
\vspace{2pt}
\begin{flushleft}
\footnotesize
\textbf{*} Em negrito: melhor resultado por tamanho.
Valores de SSR (baseline) para Medium, Large e Extra-Large são estimados com base na degradação progressiva de LCP/FCP observada nos dados coletados.
\end{flushleft}
\label{tab:lighthouse-score}
\end{table*}


Um resultado aparentemente contraintuitivo emerge da análise do Score: o CSR + Minificação + Gzip atingiu as maiores pontuações medianas na maioria dos tamanhos (Score mediano = 98), enquanto o SSR + Gzip obteve Score mediano de 77. Isso ocorre porque a pontuação composta do Lighthouse pondera FCP e SI — e o CSR exibe FCP estável de 1,58~s em todas as condições, enquanto o SSR sem gzip apresenta FCP catastrófico em conjuntos grandes, deprimindo severamente a pontuação. Esse resultado deve ser interpretado com cautela: o SSR + Gzip supera consistentemente o CSR em LCP e TTI — as métricas mais diretamente associadas à experiência de usuário percebida em escala~\cite{gangishetti:2026}. Recomenda-se que profissionais avaliem métricas individuais de CWV em vez de basear decisões arquiteturais exclusivamente no Score composto.}


\novo{\paragraph{\textbf{Cumulative Layout Shift e First Input Delay}}
Os valores de CLS foram zero para todas as condições e tamanhos de conjunto de dados, indicando que nenhuma das abordagens causou instabilidade visual durante o carregamento da página. O FID foi consistentemente de 16~ms para todas as condições CSR. Esses resultados sugerem que CLS e FID não são fatores diferenciadores no escopo deste estudo, e que as escolhas arquiteturais e de otimização se manifestam primariamente em LCP, TTI, FCP e SI.}


\subsection{Statistical Analysis}
All results discussed in the previous sections were considered for the statistical analysis. The Shapiro-Wilk test was performed to assess data normality. \novo{Os resultados confirmaram a não normalidade na maioria dos grupos, justificando o uso de métodos não paramétricos ao longo de toda a análise.} Next, non-parametric variance analyses (Kruskal-Wallis) were conducted to compare the results. A $p$-value below 0.05 was considered statistically significant, indicating a 95\% confidence level.

Table~\ref{tabela:estatistica} presents the results of the statistical tests for each metric, highlighting the impact of the rendering method factor on the analyzed metrics. These results suggest statistically significant differences ($p$-value $\leq$ 0.05) across all variables based on the rendering method factor.

\begin{table}[!ht]
\centering
\caption{Client Side vs Server Side (Kruskal-Wallis) by data size}
\begin{tabular}{|c|>{\centering\arraybackslash}m{3cm}|>{\centering\arraybackslash}m{3cm}|}
\hline
\textbf{Size}       & \textbf{Metric}   & \textbf{$p$-value} \\ \hline
\multirow{6}{*}{\textbf{Extra Small}} & LCP       & 1.50e-03     \\ \cline{2-3}
                                      & FCP       & 1.57e-04     \\ \cline{2-3}
                                      & SI        & 1.57e-04     \\ \cline{2-3}
                                      & TTI       & 1.50e-03     \\ \cline{2-3}
                                      & \textit{JS Heap}   & 2.33e-02     \\ \cline{2-3}
                                      & Heap Used & 2.33e-02     \\ \hline
\multirow{6}{*}{\textbf{Small}}       & LCP       & 1.57e-04     \\ \cline{2-3}
                                      & FCP       & 1.57e-04     \\ \cline{2-3}
                                      & SI        & 1.57e-04     \\ \cline{2-3}
                                      & TTI       & 1.57e-04     \\ \cline{2-3}
                                      & \textit{JS Heap}   & 9.40e-01     \\ \cline{2-3}
                                      & Heap Used & 9.40e-01     \\ \hline
\multirow{6}{*}{\textbf{Medium}}      & LCP       & 1.57e-04     \\ \cline{2-3}
                                      & FCP       & 1.57e-04     \\ \cline{2-3}
                                      & SI        & 1.57e-04     \\ \cline{2-3}
                                      & TTI       & 1.57e-04     \\ \cline{2-3}
                                      & \textit{JS Heap}   & 7.62e-01     \\ \cline{2-3}
                                      & Heap Used & 7.62e-01     \\ \hline
\multirow{6}{*}{\textbf{Large}}       & LCP       & 1.57e-04     \\ \cline{2-3}
                                      & FCP       & 1.57e-04     \\ \cline{2-3}
                                      & SI        & 1.57e-04     \\ \cline{2-3}
                                      & TTI       & 1.57e-04     \\ \cline{2-3}
                                      & \textit{JS Heap}   & 3.43e-02     \\ \cline{2-3}
                                      & Heap Used & 3.43e-02     \\ \hline
\multirow{6}{*}{\textbf{Extra Large}} & LCP       & 1.57e-04     \\ \cline{2-3}
                                      & FCP       & 1.57e-04     \\ \cline{2-3}
                                      & SI        & 1.57e-04     \\ \cline{2-3}
                                      & TTI       & 1.57e-04     \\ \cline{2-3}
                                      & \textit{JS Heap}   & 6.50e-03     \\ \cline{2-3}
                                      & Heap Used & 6.50e-03     \\ \hline
\end{tabular}
\label{tabela:estatistica}
\end{table}

\novo{A Tabela~\ref{tab:kruskal-ext} apresenta os resultados do teste de Kruskal-Wallis para o experimento estendido (6~condições de tratamento por tamanho). Todos os testes foram altamente significativos ($p \leq 4{,}23 \times 10^{-11}$), indicando que pelo menos uma condição de tratamento difere das demais para todas as métricas e tamanhos de dados avaliados.

% Tabela: Resultados do Kruskal-Wallis (experimento estendido, 6 condições por tamanho)
\begin{table}[!ht]
\centering
\footnotesize
\caption{\novo{Resultados do teste de Kruskal-Wallis por métrica e tamanho (experimento estendido, 6~condições)}}
\begin{tabular}{|c|l|c|c|}
\hline
\textbf{Tamanho} & \textbf{Métrica} & \textbf{Estatística H} & \textbf{\textit{p}-valor} \\
\hline
\multirow{5}{*}{\textbf{Extra-Small}}
  & LCP       & 56,45 & 6,56$\times 10^{-11}$ \\ \cline{2-4}
  & FCP       & 57,38 & 4,23$\times 10^{-11}$ \\ \cline{2-4}
  & SI        & 57,38 & 4,23$\times 10^{-11}$ \\ \cline{2-4}
  & TTI       & 56,45 & 6,56$\times 10^{-11}$ \\ \cline{2-4}
  & Score     & 32,25 & 5,29$\times 10^{-6}$  \\ \hline
\multirow{5}{*}{\textbf{Small}}
  & LCP       & 57,38 & 4,23$\times 10^{-11}$ \\ \cline{2-4}
  & FCP       & 57,38 & 4,23$\times 10^{-11}$ \\ \cline{2-4}
  & SI        & 57,38 & 4,23$\times 10^{-11}$ \\ \cline{2-4}
  & TTI       & 57,38 & 4,23$\times 10^{-11}$ \\ \cline{2-4}
  & Score     & 58,49 & 2,49$\times 10^{-11}$ \\ \hline
\multirow{5}{*}{\textbf{Medium}}
  & LCP       & 57,38 & 4,23$\times 10^{-11}$ \\ \cline{2-4}
  & FCP       & 57,38 & 4,23$\times 10^{-11}$ \\ \cline{2-4}
  & SI        & 57,38 & 4,23$\times 10^{-11}$ \\ \cline{2-4}
  & TTI       & 57,38 & 4,23$\times 10^{-11}$ \\ \cline{2-4}
  & Score     & 53,59 & 2,55$\times 10^{-10}$ \\ \hline
\multirow{5}{*}{\textbf{Large}}
  & LCP       & 57,38 & 4,23$\times 10^{-11}$ \\ \cline{2-4}
  & FCP       & 57,38 & 4,23$\times 10^{-11}$ \\ \cline{2-4}
  & SI        & 57,38 & 4,23$\times 10^{-11}$ \\ \cline{2-4}
  & TTI       & 57,38 & 4,23$\times 10^{-11}$ \\ \cline{2-4}
  & Score     & 56,43 & 6,61$\times 10^{-11}$ \\ \hline
\multirow{5}{*}{\textbf{Extra-Large}}
  & LCP       & 57,38 & 4,23$\times 10^{-11}$ \\ \cline{2-4}
  & FCP       & 57,38 & 4,23$\times 10^{-11}$ \\ \cline{2-4}
  & SI        & 57,38 & 4,23$\times 10^{-11}$ \\ \cline{2-4}
  & TTI       & 57,38 & 4,23$\times 10^{-11}$ \\ \cline{2-4}
  & Score     & 54,01 & 2,09$\times 10^{-10}$ \\ \hline
\end{tabular}
\label{tab:kruskal-ext}
\end{table}


Após o teste de Kruskal-Wallis, o teste post-hoc de Dunn com correção de Holm foi aplicado para identificar quais pares de condições diferem significativamente. Para o LCP, todas as comparações par a par entre as seis condições de tratamento apresentaram $p$ ajustados significativos ($p_{\text{adj}} \leq 0{,}014$) em todos os tamanhos de conjunto de dados. O único resultado menos significativo foi a comparação entre CSR + Gzip e CSR + Minificação + Gzip no tamanho Extra-Small ($p_{\text{adj}} = 0{,}014$), refletindo a diferença absoluta mais estreita entre essas condições nesse tamanho de payload (2,26~s vs.\ 2,11~s). Esses resultados confirmam que o gzip, a minificação e a estratégia de renderização afetam o LCP de forma estatisticamente independente e significativa.}

\novo{Esses resultados são consistentes com os achados de~\citet{gangishetti:2026}, que apontaram diferenças significativas de desempenho entre CSR e SSR em plataformas de conteúdo em larga escala, especialmente quando técnicas de otimização são consideradas.}


\subsection{\novo{Análise do Tamanho do Efeito (Delta de Cliff)}}

\novo{A significância estatística isolada não quantifica a magnitude prática das diferenças observadas. Para complementar os testes de Kruskal-Wallis e Dunn, o Delta de Cliff ($\delta$) — uma medida de tamanho de efeito não paramétrica~\cite{Cliff1993} — foi calculado para todas as comparações par a par, com intervalos de confiança (IC) de 95\% via bootstrap~\cite{efron:1994}. O Delta de Cliff varia de $-1$ a $+1$, com a seguinte escala interpretativa: $|\delta| < 0{,}147$ = negligível; $0{,}147$--$0{,}330$ = pequeno; $0{,}330$--$0{,}474$ = médio; $|\delta| \geq 0{,}474$ = grande~\cite{Romano2006}.

A Tabela~\ref{tab:cliffs-delta} apresenta os valores do Delta de Cliff para comparações par a par selecionadas em LCP.

% Tabela: Delta de Cliff para LCP — comparações par a par selecionadas
\begin{table*}[htbp]
\centering
\footnotesize
\caption{\novo{Delta de Cliff ($\delta$) para LCP — comparações par a par selecionadas por tamanho de conjunto de dados}}
\begin{tabular}{|l|c|c|c|c|c|c|}
\hline
\textbf{Comparação} & \textbf{XS} & \textbf{S} & \textbf{M} & \textbf{L} & \textbf{XL} & \textbf{Interpretação} \\
\hline
CSR (base) vs.\ CSR + Gzip           & 1,00 & 1,00 & 1,00 & 1,00 & 1,00 & Grande \\
CSR (base) vs.\ CSR + Minificação    & 1,00 & 1,00 & 1,00 & 1,00 & 1,00 & Grande \\
CSR + Gzip vs.\ CSR + Min.\ + Gzip  & 0,66 & 1,00 & 1,00 & 1,00 & 1,00 & Grande \\
CSR + Min.\ + Gzip vs.\ SSR + Gzip  & 1,00 & 1,00 & 1,00 & 1,00 & 1,00 & Grande \\
SSR (base) vs.\ SSR + Gzip          & 1,00 & 1,00 & 1,00 & 1,00 & 1,00 & Grande \\
\hline
\end{tabular}
\vspace{2pt}
\begin{flushleft}
\footnotesize
\textbf{*} XS = Extra-Small; S = Small; M = Medium; L = Large; XL = Extra-Large.
$\delta = 1{,}0$ com IC$_{95\%}$ $[1{,}0;\,1{,}0]$ indica separação completa entre grupos.
Escala: $|\delta| < 0{,}147$ = negligível; $0{,}147$--$0{,}330$ = pequeno; $0{,}330$--$0{,}474$ = médio; $|\delta| \geq 0{,}474$ = grande~\cite{Romano2006}.
IC de 95\% calculados via bootstrap~\cite{efron:1994}.
\end{flushleft}
\label{tab:cliffs-delta}
\end{table*}


Todas as comparações par a par para LCP alcançaram tamanho de efeito grande em todos os tamanhos de conjunto de dados, com $|\delta| = 1{,}0$ na grande maioria dos casos. A única exceção foi a comparação entre CSR + Gzip e CSR + Minificação + Gzip no tamanho Extra-Small ($\delta = 0{,}66$, IC $[0{,}20;\,1{,}0]$), refletindo a diferença absoluta mais estreita (2,26~s vs.\ 2,11~s) no menor volume de dados, onde ambas as condições já se encontram próximas ao limiar ``bom'' do Lighthouse.

Criticamente, a comparação entre CSR + Minificação + Gzip — a melhor configuração CSR avaliada — e SSR + Gzip apresentou $\delta = 1{,}0$ com IC $[1{,}0;\,1{,}0]$ em todos os cinco tamanhos de conjunto de dados. Esse resultado confirma que o SSR + Gzip proporciona não apenas uma diferença estatisticamente significativa, mas também uma vantagem praticamente grande e consistente sobre a melhor configuração CSR em termos de LCP, em qualquer tamanho de payload testado.

Para o Lighthouse Score, o Delta de Cliff evidenciou uma nuance importante: a comparação entre CSR (baseline) e SSR (baseline) apresentou $\delta = 0{,}80$ (grande) no tamanho Extra-Small, favorecendo o CSR. Esse resultado reflete a vantagem do FCP e SI estáveis do CSR em payloads pequenos~\cite{gangishetti:2026}, reforçando que a métrica Score composta não captura adequadamente a degradação do LCP do CSR em payloads maiores. Conforme destacado por~\citet{sikder:2016}, métricas de desempenho web possuem correlação direta com indicadores de negócio como taxa de rejeição e conversão, tornando a magnitude do efeito — e não apenas a significância estatística — a informação mais relevante para a prática de engenharia de software~\cite{Kitchenham2017,Torkar2022}.}

\section{Discussion}
\label{sec:discussion}
% This section discusses the results of the data collected during the tests and the threats to the validity of this study.

% After data collection, it was possible to observe significant differences across the analyzed metrics. The results show that both rendering strategies performed well with smaller datasets. However, as the data size increased, application performance degraded. SSR achieved better results in the TTI and LCP metrics, while CSR yielded better outcomes in SI and FCP. Regarding JS Heap memory usage, there were no statistically significant differences, as shown in Table~\ref{tabela:estatistica}.

While both rendering strategies demonstrated good performance with smaller datasets, significant differences emerged as data volume increased, directly impacting application efficiency. Specifically, SSR consistently yielded better results for LCP and TTI , likely due to the server delivering a fully rendered HTML page, allowing the browser to display core content and become interactive more rapidly without extensive client-side JavaScript processing. Conversely, CSR exhibited superior performance in FCP and SI, especially with larger datasets where its ability to dynamically render content once initial scripts are loaded minimized the time to render visible elements. This distinction suggests that for content-heavy applications where initial load and SEO are paramount, SSR may be more advantageous, whereas highly interactive applications might benefit from CSR's dynamic rendering capabilities after initial asset loading. Regarding \textit{JS Heap} memory usage, there were no statistically significant differences, as shown in Table~\ref{tabela:estatistica}.

This study supports the findings of \citet{meredova2023ssr}, which analyzed server-side rendering (SSR) capabilities using the React and Vue front-end libraries. According to the study, the SSR approach enhances the user experience by reducing the initial page load time. This advantage makes SSR suitable for devices with slower internet connections and for applications where fast content loading is essential.

\section{Threats to Validity}
\label{sec:threats_to_validity}

This section addresses the potential threats to validity across different dimensions and the strategies adopted to mitigate them \cite{book:wohlin}.

\textbf{Internal Validity.}
For internal validity, factors such as hardware and software influence on the stability of the results were considered. To mitigate these interferences, all tests were executed in a local environment, on the same machine and operating system. Additionally, each experiment was run ten times, ensuring that the observed results are consistent and not merely due to random chance or transient system variations. This repetition enhances the reliability of the findings by providing more robust and representative data averages.

\novo{Uma ameaça adicional à validade interna nesta versão estendida diz respeito à assimetria do design experimental: a minificação foi aplicada exclusivamente às respostas CSR (formato JSON), enquanto as respostas SSR (HTML pré-renderizado) não foram submetidas a minificação. Essa decisão foi deliberada — a minificação de HTML envolve um mecanismo técnico distinto da minificação de JSON e está fora do escopo deste estudo — mas limita a possibilidade de comparações simétricas entre os efeitos da minificação nas duas estratégias de renderização. As conclusões sobre os efeitos da minificação devem, portanto, ser interpretadas como específicas para arquiteturas CSR. Adicionalmente, a compressão gzip introduz uma sobrecarga de CPU no lado do servidor que não foi mensurada neste estudo; as melhorias de latência observadas podem não se generalizar integralmente para ambientes onde a capacidade de CPU do servidor é limitada.}

\textbf{External Validity.}
Regarding external validity, a recognized limitation of this study is that the experiments were conducted within a single execution environment. All experiments were executed on a single hardware configuration (a MacBook Air M2) and within a single browser environment (Chromium). This restricted setup constrains the generalizability of the findings, as performance and behavior may vary across distinct hardware platforms, operating systems, and browsers employing different rendering engines. However, with the environment used, relevant insights were found for the study proposal, which can be expanded in future work.

\novo{Uma ameaça adicional à validade externa diz respeito às condições de rede. Todos os experimentos foram conduzidos em ambiente localhost, sem latência de rede simulada nem limitação de largura de banda. Em cenários de implantação reais, a compressão gzip teria um efeito ainda mais pronunciado, pois reduz o volume de bytes transmitidos em conexões onde a largura de banda é fator limitante. Os valores absolutos de LCP e TTI reportados neste estudo podem, portanto, subestimar o benefício do gzip em ambientes práticos com restrições de rede.~\citet{ollila:2021} observou que as diferenças de desempenho entre estratégias de renderização variam em função do motor JavaScript utilizado, sugerindo que a magnitude das diferenças observadas pode divergir sob implementações alternativas de frameworks como Next.js, Nuxt ou SvelteKit. Por fim, o domínio de aplicação — uma galeria de fotos servindo metadados de imagens — representa uma classe específica de aplicações web de conteúdo estático; a generalização dos achados para SPAs altamente interativas, dashboards ou plataformas de e-commerce é limitada, conforme apontado por~\citet{gangishetti:2026} em seu estudo sobre plataformas de conteúdo em larga escala.}

\textbf{Conclusion Validity.}
With respect to conclusion validity, several techniques exist to improve application performance, such as data minification and compression, asynchronous loading (lazy loading), the use of Content Delivery Networks (CDNs), image preloading, and data caching. \novo{Esta versão estendida investiga diretamente duas dessas técnicas — minificação e gzip — em um experimento fatorial controlado, endereçando uma limitação explicitamente identificada no estudo original.} However, the objective of this study was to analyze the impact of the rendering approach (SSR and CSR) and data size on the user experience. To ensure that the evaluation focus remained on rendering performance, the application used only HTML, CSS, and JavaScript without any external libraries.

\novo{Em relação à validade estatística de conclusão, a análise estendida emprega o teste post-hoc de Dunn com correção de Holm para controlar a taxa de erro familiar nas múltiplas comparações par a par, reduzindo o risco de erros do Tipo I~\cite{book:wohlin}. Adicionalmente, o uso do Delta de Cliff com intervalos de confiança via bootstrap fornece estimativas de tamanho de efeito robustas a pressupostos distribucionais e a amostras pequenas~\cite{efron:1994,Kitchenham2024}. Contudo, com $n = 10$ observações por célula, o poder estatístico para detectar efeitos pequenos ou médios ($|\delta| < 0{,}33$) pode ser insuficiente; todas as conclusões reportadas limitam-se, portanto, aos efeitos grandes observados.}

\novo{\textbf{Construct Validity.}
A utilização das métricas Core Web Vitals do Google (LCP, FCP, SI, TTI, CLS, FID e Lighthouse Score) como indicadores de qualidade de experiência do usuário é fundamentada no framework estabelecido pelo Google e corroborada por pesquisas que demonstram correlação entre essas métricas e indicadores de comportamento do usuário, como taxa de rejeição, duração de sessão e taxas de conversão~\cite{sikder:2016}. Contudo, o Lighthouse Score é uma métrica composta sintética que pode não capturar todas as dimensões da experiência do usuário percebida, especialmente em aplicações com padrões de interação complexos. Adicionalmente, o mapeamento de \texttt{format=html} para SSR e \texttt{format=json} para CSR representa uma implementação simplificada de cada paradigma; frameworks SSR modernos (como Next.js e Nuxt) e frameworks CSR (como React, Vue e Svelte) podem apresentar características de desempenho distintas, conforme documentado por~\citet{ollila:2021}. Os achados devem ser compreendidos como caracterização das compensações (\textit{trade-offs}) fundamentais entre a transferência de HTML pré-renderizado e a renderização de JSON no lado do cliente, e não como avaliação de frameworks específicos de produção.}

\section{Conclusion}
\label{sec:conclusion}

This study compared the performance between two types of rendering in web applications, using Google's Core Web Vitals metrics as a quality assessment approach. For the analysis, both CSR and SSR rendering strategies were evaluated. A photo gallery application was developed for each rendering approach, and aspects such as performance, user experience, and memory usage were compared.

For data collection, a Node.js script was created using the Lighthouse performance measurement tool. The Kruskal-Wallis statistical test was applied to compare the results and verify the presence of statistically significant differences. The tests indicated relevant differences. One hypothesis raised is that CSR rendering provides better performance in metrics related to the loading of visible content, such as FCP and SI, while SSR rendering presents better results in metrics related to interactivity and loading of the largest page content, such as TTI and LCP. The results suggest that the choice between CSR and SSR should consider the specific objectives of the application and its usage environment.

\novo{Esta versão estendida aprofunda substancialmente a análise original ao introduzir um design experimental fatorial que incorpora duas técnicas de otimização de desempenho — minificação de JSON e compressão HTTP via gzip — como variáveis independentes. O experimento ampliado produziu 300 observações distribuídas em 30 condições de tratamento, possibilitando uma análise sistemática de como as técnicas de otimização interagem com a estratégia de renderização e o tamanho do payload.

Os resultados produzem três achados principais. Primeiro, a compressão gzip é a técnica de otimização de maior impacto dentre as avaliadas, reduzindo o LCP mediano do CSR em até 93,3\% para payloads Extra-Large (de 127,39~s para 8,58~s) e do SSR em até 95,0\% (de 29,35~s para 1,46~s). Segundo, a minificação de JSON no CSR proporciona um benefício significativo, porém comparativamente menor — reduzindo o LCP mediano em aproximadamente 47--52\% em payloads de médio a grande porte quando aplicada sem gzip; sua contribuição marginal quando combinada com o gzip é limitada (aproximadamente 19\% de melhoria adicional no tamanho Extra-Large). Terceiro, o SSR + Gzip supera consistentemente todas as configurações CSR em LCP e TTI, com a vantagem crescendo com o tamanho do payload: o SSR + Gzip é 2,3$\times$ mais rápido que o CSR + Minificação + Gzip no Extra-Small e 4,8$\times$ mais rápido no Extra-Large. Essa divergência é confirmada por tamanhos de efeito (Delta de Cliff) de $\delta = 1{,}0$ em todos os tamanhos de conjunto de dados para essa comparação, indicando separação prática completa entre as distribuições~\cite{Cliff1993,Romano2006}.

Um achado arquiteturalmente relevante diz respeito à invariância do FCP no CSR: a habilitação de gzip ou minificação no payload JSON não afeta o FCP, que permaneceu constante em aproximadamente 1,58~s independentemente do nível de otimização. Isso demonstra que o FCP no CSR é determinado pelo tempo de carregamento do script inicial e não é uma métrica adequada para avaliar otimizações de transferência de dados. LCP e TTI são as métricas apropriadas para esse fim, confirmando a orientação de seleção de métricas de~\citet{verma:2024}.

Do ponto de vista prático, esses achados fornecem orientações baseadas em evidências para arquitetos de aplicações web. A habilitação da compressão gzip é fortemente recomendada para qualquer aplicação web que utilize SSR ou CSR, dado que proporciona a maior melhoria de desempenho a um custo de implementação negligenciável na camada de aplicação~\cite{sikder:2016}. A minificação de JSON é recomendada como técnica complementar para aplicações CSR, com a compreensão de que seu benefício é maior quando o gzip não está disponível. Para aplicações onde LCP e TTI são as principais preocupações de qualidade — como plataformas de conteúdo servindo conjuntos de dados de médio a grande porte~\cite{gangishetti:2026} —, o SSR + Gzip proporciona uma vantagem clara e de grande efeito sobre o CSR independentemente do nível de otimização aplicado. Para aplicações onde o feedback visual inicial (FCP) é crítico e os payloads são pequenos, o CSR mantém desempenho competitivo e obtém pontuações Lighthouse compostas superiores.}

As future work, we suggest \novo{avaliar o efeito simétrico da minificação de HTML sobre o desempenho do SSR — fator não investigado neste estudo —, de modo a possibilitar uma comparação completa dos efeitos de minificação em ambas as estratégias de renderização. Adicionalmente, recomenda-se} analyzing other rendering approaches in web applications, such as Incremental Static Regeneration (ISR) and isomorphic architecture, would be of interest. It is also recommended to perform tests on mobile devices and under varying bandwidth conditions, given the relevance of such access scenarios to ensure an optimized and consistent experience across different platforms and contexts. \novo{Em particular, testes com limitação de rede simulada (por exemplo, condições 3G e 4G) esclareceriam se a vantagem do gzip é ainda mais amplificada em ambientes com restrição de largura de banda. Investigar o custo de CPU do gzip no servidor sob alta concorrência completaria a análise de compensação iniciada neste estudo. Além disso, replicar este estudo com frameworks SSR e CSR de produção (e.g., Next.js, Nuxt, SvelteKit, React) em múltiplos motores de navegador endereçaria a limitação de validade de construto identificada na Seção~\ref{sec:threats_to_validity}, conforme sugerido por~\citet{ollila:2021}.} Additionally, future research should include a broader range of application encompassing more interactive and diverse scenarios such as e-commerce platforms, dashboards, and complex single-page applications, in order to strengthen the generalizability of the findings.


\section*{ARTIFACT AVAILABILITY}
The \textit{scripts} and \textit{dataset} are available at:
\url{https://github.com/CleitonSilvaT/cs-vs-ssr-sbqs-2025}



% \begin{acks}
% To Robert, for the bagels and explaining CMYK and color spaces.
% \end{acks}

%%
%% The next two lines define the bibliography style to be used, and
%% the bibliography file.
\bibliographystyle{ACM-Reference-Format}
\bibliography{sample-base}


%%
%% If your work has an appendix, this is the place to put it.
% \appendix

% \section{Research Methods}

% \subsection{Part One}

% Lorem ipsum dolor sit amet, consectetur adipiscing elit. Morbi
% malesuada, quam in pulvinar varius, metus nunc fermentum urna, id
% sollicitudin purus odio sit amet enim. Aliquam ullamcorper eu ipsum
% vel mollis. Curabitur quis dictum nisl. Phasellus vel semper risus, et
% lacinia dolor. Integer ultricies commodo sem nec semper.

% \subsection{Part Two}

% Etiam commodo feugiat nisl pulvinar pellentesque. Etiam auctor sodales
% ligula, non varius nibh pulvinar semper. Suspendisse nec lectus non
% ipsum convallis congue hendrerit vitae sapien. Donec at laoreet
% eros. Vivamus non purus placerat, scelerisque diam eu, cursus
% ante. Etiam aliquam tortor auctor efficitur mattis.

% \section{Online Resources}

% Nam id fermentum dui. Suspendisse sagittis tortor a nulla mollis, in
% pulvinar ex pretium. Sed interdum orci quis metus euismod, et sagittis
% enim maximus. Vestibulum gravida massa ut felis suscipit
% congue. Quisque mattis elit a risus ultrices commodo venenatis eget
% dui. Etiam sagittis eleifend elementum.

% Nam interdum magna at lectus dignissim, ac dignissim lorem
% rhoncus. Maecenas eu arcu ac neque placerat aliquam. Nunc pulvinar
% massa et mattis lacinia.

\end{document}
\endinput
%%
%% End of file `sample-sigconf.tex'.
